\begin{table}[!htbp]\centering\scriptsize\setlength{\tabcolsep}{3pt}
\caption{Architectures comparées, à recette égale (220~Md€)}\label{tab:ircsg_arch}
\resizebox{\linewidth}{!}{%
\begin{tabular}{@{}lrrrrrrr@{}}
\toprule
& Points & Gain d'efficacité & Part du & Perdants & Retraités & Couple, & 10~\% \\
& pour tous & Md€ par an & gain & $>$ 500~€ & Md€ & 2 enfants & du haut \\
& & (fourchette) & atteignable & & & € par foyer & Md€ \\
\midrule
Barème neutre & 0,0 & 4,0 (1,9 à 9,2) & 52~\% & 25~\% & +7,0 & +570 & +2,5 \\
\multicolumn{8}{@{}l}{\emph{Un dispositif ajouté au barème neutre}} \\
\quad Franchise d'activité 4~000~€ & +1,8 & 4,6 (2,1 à 10,2) & 59~\% & 30~\% & +14,0 & +217 & +3,6 \\
\quad Franchise d'activité 8~000~€ & +3,6 & 5,0 (2,2 à 11,0) & 64~\% & 33~\% & +21,4 & $-$172 & +4,7 \\
\quad Franchise de 8~000~€ sur tous les revenus & +5,8 & 3,3 (1,4 à 7,4) & 43~\% & 21~\% & +5,4 & +775 & +9,6 \\
\quad Franchise de 6~000~€ par enfant & +1,5 & 3,6 (1,7 à 8,4) & 47~\% & 27~\% & +13,0 & $-$983 & +0,6 \\
\quad Chèque de 742~€ par enfant & +0,9 & 3,1 (1,5 à 7,3) & 40~\% & 26~\% & +10,7 & $-$361 & +3,1 \\
\quad Quotient conjugal et familial & +2,0 & 3,9 (1,9 à 9,0) & 50~\% & 28~\% & +12,1 & $-$965 & $-$3,0 \\
\quad Abattement de 10~\% sur les pensions & +0,6 & 3,4 (1,6 à 8,0) & 44~\% & 24~\% & +0,9 & +923 & +2,7 \\
\quad Petites pensions exonérées & +0,4 & 3,7 (1,7 à 8,5) & 47~\% & 16~\% & +3,2 & +786 & +4,0 \\
\midrule
\textbf{Proposition}, franchise d'activité 4~000~€ & --- & 4,9 (2,3 à 11,1) & 63~\% & 29~\% & +15,2 & $-$1\,195 & $-$1,7 \\
Proposition, franchise d'activité 8~000~€ & --- & 5,2 (2,4 à 11,8) & 67~\% & 26~\% & +22,8 & $-$1\,629 & $-$1,5 \\
\midrule
Optimum d'efficacité, sans franchise & --- & 7,4 (3,7 à 17,4) & 95~\% & 58~\% & +13,7 & $-$519 & $-$25,7 \\
Gain atteignable (optimum, franchise 4~000~€) & --- & 7,8 (3,8 à 18,2) & 100~\% & 49~\% & +20,3 & $-$869 & $-$25,1 \\
\bottomrule
\end{tabular}}
\par\vspace{3pt}\parbox{\linewidth}{\raggedright\scriptsize\color{encre}
Gain d'efficacité : baisse de la perte sèche sur l'activité réelle par rapport à aujourd'hui, hypothèses centrales (basses et hautes entre parenthèses) ; la réaction du revenu déclaré n'y est pas comptée. Points pour tous : hausse des taux du barème neutre qui finance le dispositif. Proposition : franchise d'activité, 6~000~€ par enfant, pensions exonérées jusqu'à 12~500~€ puis de moins en moins jusqu'à 25~000~€ ; barème le plus efficace parmi ceux qui ne changent la charge moyenne d'aucun décile de plus de 500~€ et n'alourdissent pas celle du 1~\% le plus aisé. Retraités, couples, 10~\% du haut : variation de leur impôt (positif : ils paient plus).}
\end{table}
